% Copy-paste in Paper (Pfad A, nach QQ-Spektraltabelle)
\begin{definition}[Arithmetischer Resonanzpunkt]
Sei $K=\mathbb{Q}(\zeta_5)$, $V_{\mathrm{arith}}\supset \mathbb{Q}^4$ ein $K$-Modul mit Rundweg-Operator
$\Pi_\Gamma=\frac14(I+\Gamma+\Gamma^2+\Gamma^3)$. Ein $x\in V_{\mathrm{arith}}$ heißt
\emph{arithmetischer Resonanzpunkt}, wenn $\Pi_\Gamma x=x$ und $\sigma(x)\neq x$ für ein
$\sigma\in\mathrm{Gal}(K/\mathbb{Q})$ mit $\sigma(\Pi_\Gamma)=\Pi_\Gamma$.
\end{definition}
\begin{remark}
Auf dem reinen Lift $K^4$ mit $\mathbb{Q}$-Eigenbasen ist $\sigma|_{\mathbb{Q}^4}=\mathrm{id}$;
der erste nichttriviale Test identifiziert $(E,A,B,C)$ mit $\zeta,\ldots,\zeta^4$ und prüft
$\varphi_p$ für $p\neq 5$ auf gemischten Vektoren in $\ker(\Pi_\Gamma-I)$.
\end{remark}
